% 11-08
\section{Introduction}
\label{sec:intro}

Natural conversation requires both speaking and listening. While one
participant holds the conversational floor, the listener conveys
attention and engagement through facial expressions, head movements,
and brief vocal acknowledgments. These backchannels support the
ongoing utterance without taking over the conversation. Visual
reactions may occur alone or accompany vocal feedback, while quiet
listening is also appropriate. For digital humans, natural interaction
therefore depends not only on producing fluent spoken responses, but
also on how they behave while the user is speaking. Realistic faces
and voices alone do not establish whether these listening behaviors
are appropriate to the conversational context.

Recent advances have substantially expanded the capabilities of
digital human generation. Speech-driven talking-head methods, such
as OmniHuman-1.5~\cite{omnihuman15}, synthesize visual behavior from
the target person's supplied speech. Dedicated listening-generation
methods, including ViCo~\cite{vico}, Learning to
Listen~\cite{learningtolisten}, and GLARE~\cite{glare}, instead
generate reactions to another speaker. GLARE further introduces
reaction-oriented supervision and evaluation, but these methods
focus on visual outputs without generating the listener's own
vocal feedback. In settings that supply listener audio, part of
the response is already prescribed, leaving the generation of
that audio to another component. OmniResponse~\cite{omniresponse}
and OmniPersona~\cite{omnipersona} advance beyond these settings
by generating audiovisual conversational responses, including
listener feedback. Their broader conversational scope motivates
a complementary question: how can learning specifically focused
on the listening phase improve the audiovisual feedback expressed
within it? We treat listening feedback as the primary learning
and evaluation target, rather than assessing it only as part of
a complete conversational exchange.

We study audiovisual backchannel generation from the current
speaker audio and a reference image of the listener
(\cref{fig:task}). This setting supports interaction in which a
user speaks to a visible digital avatar without enabling a
camera. It also accommodates initial exchanges and standalone
interactions for which no prior dialogue history is available.
The reference image specifies the avatar's visual identity,
while the current utterance provides the conversational input.
The listener's own audio is an output to be generated, rather
than a supplied condition that prescribes whether, when, and
what the avatar vocalizes. Our goal is to generate visual
reactions, brief vocal feedback, and periods of quiet listening
while the speaker retains the floor.

Learning this behavior presents two related challenges. First,
the semantic and acoustic information in the speaker's speech
must be translated into conditions suitable for the listener's
audiovisual response. Unlike speech-driven animation, the input
audio belongs to the other participant and does not specify
the target person's articulation or intended reaction. The
conditioning representations must therefore be adapted to
listener generation. Second, the training objective must
account for feedback behaviors that can be brief, subtle, and
interspersed with quiet listening. Media reconstruction
supervises the audiovisual output as a whole, but does not
separately specify which feedback categories should be
expressed. We investigate explicit event supervision as a
complement to reconstruction, with the aim of making these
behaviors a direct target of learning.

To support this learning, we construct a listening-focused
audiovisual dataset from Seamless
conversations~\cite{seamless}. We organize paired clips in
which the speaker retains the conversational floor, excluding
segments where the listener takes over and develops a spoken
response while preserving brief vocal backchannels. This
distinction separates listening feedback from turn-taking
responses without equating listening with silence. Our
annotations cover both visual behavior and the listener's
own sounds. We provide extracted visual and vocal cues,
together with their temporal intervals, to
Qwen3-Omni~\cite{qwen3omni} to support detailed descriptions
of observable reactions. The descriptions retain subtle,
brief, and compound behaviors and organize them in temporal
order, rather than reducing each clip to a coarse summary.
The resulting samples associate speaker audio with listener
audiovisual media and fine-grained feedback descriptions.

Our generation framework combines direct speaker-audio
conditioning with supervision of both conditioning
representations and feedback events. Speaker-audio
cross-attention supplies acoustic information to the
generator. During training, features extracted from the
target listener audiovisual media supervise the conditioning
representations, helping adapt them to listener generation.
Clip-level event supervision additionally encourages the
expression of the corresponding feedback categories alongside
media reconstruction. These designs address how speaker
information reaches the generator and how feedback behavior
is represented in the learning objective. The event
supervision operates at the clip level, allowing feedback
timing to vary rather than imposing exact event timestamps
through this objective.

We evaluate generation specifically during listening,
considering audiovisual quality, feedback-category agreement,
contextual appropriateness and engagement, and audiovisual
coordination. A speaker utterance may admit multiple
reasonable reactions, so matching a single recorded response
cannot fully determine feedback quality. Clip-level category
agreement measures whether reference feedback categories are
expressed without requiring exact temporal alignment.
Contextual evaluation additionally considers whether
alternative responses fit the utterance and preserve the
listening role, while audiovisual coordination concerns the
relationship between generated sounds and movements. This
evaluation distinguishes the quality of the synthesized
media from the listening behavior it conveys.

Our contributions are as follows:
\begin{itemize}
    \item We introduce an audiovisual backchannel generation
    framework conditioned on current speaker audio and a
    listener reference image. It combines direct acoustic
    conditioning, target audiovisual feature supervision,
    and clip-level event supervision to support feedback
    generation during listening.

    \item We construct listening-focused audiovisual data
    through conversational-role filtering and fine-grained
    feedback annotation. The data distinguish backchannels
    from turn-taking responses and associate speaker audio
    with listener audiovisual media and detailed descriptions
    of visual and vocal behavior.

    \item We formulate a listening-focused evaluation design
    that complements audiovisual quality with feedback-category
    agreement, contextual appropriateness and engagement,
    and audiovisual coordination, accommodating reasonable
    variation in listener responses.
\end{itemize}

% 11-07
% \section{Introduction}
% \label{sec:intro}
% % Paragraph role P01: 开场：任务及目标
% Speakers and listeners both contribute to natural conversation. Listeners use nods, facial expressions, and brief vocal acknowledgments to convey attention and engagement. Digital humans therefore need to produce natural audiovisual feedback while listening. We study audiovisual backchannel generation from the current speaker audio and a reference image of the listener, with the aim of making conversation more natural and engaging (\cref{fig:task}).

% \begin{figure*}[t]
% \centering
% \includegraphics[width=\textwidth]{F1_task.pdf}
% \caption{\figdraft{F1: Audiovisual backchannels during listening} Given current speaker audio and the listener's first frame, the task is to generate listener video and audio while the speaker holds the floor. Quiet listening is also valid. The rows illustrate visual and audiovisual output scopes, not measured comparisons between methods. Portraits, utterances, and waveforms are schematic.}
% \label{fig:task}
% \end{figure*}

% % Paragraph role P02: 已有能力与本文范围
% Speech-driven talking-head methods, such as OmniHuman-1.5~\cite{omnihuman15}, generate facial motion from the target person's given speech, emphasizing lip synchronization and natural motion. Listening-generation methods instead model responses to a speaker. ViCo~\cite{vico}, Learning to Listen~\cite{learningtolisten}, and GLARE~\cite{glare} generate visual listening behavior without generating the listener's own audio. More recent work, including OmniResponse~\cite{omniresponse} and OmniPersona~\cite{omnipersona}, also studies vocal and silent responses. We focus on feedback produced while the speaker holds the conversational floor, using only the current speaker audio and the listener's reference image, without dialogue history or a supplied listener audio track.

% % Paragraph role P03: 技术困难与设计目标
% This task requires converting semantic and acoustic information in the current speech into listener sounds and visual behavior. General multimodal representations provide a basis for understanding, but their mapping to audiovisual generation conditions still needs to be adapted to listening. Moreover, media reconstruction objectives do not directly specify whether feedback categories, such as nods or acknowledgments, should occur. We therefore address the use of current speech information, the adaptation of conditioning representations, and explicit supervision of feedback behavior.

% % Paragraph role P04: 数据需求
% The task also requires paired speaker inputs, listener video, and the listener's own audio, with clear listening roles and feedback annotations. Existing resources support these requirements to different extents. Some do not provide an independent listener audio track ready for generation training. Others contain both participants' audiovisual information but still require listening-oriented sample organization and backchannel annotation. Converting raw conversations into suitable training samples thus requires task-specific data preparation.
% \cnnote{待补引用}{上述资源差异沿用中文稿的限定表述。定稿时须为各类局限补充对应资源证据；不将Seamless或ResponseNet写成完全没有音轨或行为标注。}

% % Paragraph role P05: 评价需求
% Evaluation must also address listening behavior. Standard video-quality metrics describe visual fidelity and realism, but a visually convincing output may still lack meaningful feedback or respond inappropriately to the current utterance. Evaluation should therefore consider audiovisual quality, feedback behavior, naturalness, and conversational engagement.

% % Paragraph role P06: 数据方案
% We address these requirements through data construction, a generation framework, and an evaluation design. Starting from Seamless conversations~\cite{seamless}, we use pauses, ASR punctuation, and transcript-based utterance-boundary cues to organize paired speaker--listener clips. To address reaction omissions observed when describing audiovisual clips directly, we extract visual and vocal cues and provide their categories and temporal intervals to Qwen3-Omni~\cite{qwen3omni}. The resulting natural-language descriptions organize reactions in their observed order, providing paired media and descriptions for learning audiovisual listening responses.

% % Paragraph role P07: 方法方案
% Our generator uses speaker-audio cross-attention to supply direct acoustic information, target audiovisual feature supervision to adapt the conditioning representations, and clip-level event supervision to encourage the corresponding backchannel categories. These components address input information, condition adaptation, and feedback expression, respectively.
% \cnnote{证据待补}{这里陈述设计目的；三个组件的实际效果由最终同条件消融支持，不预写显著提升或最优结论。}

% % Paragraph role P08: 评价方案
% Our evaluation design considers clip-level event-category agreement, contextual appropriateness and engagement, and audiovisual event coordination. Category agreement allows feedback timing to vary, contextual evaluation allows different reasonable responses, and coordination concerns the relation between generated sounds and movements.
% \cnnote{协议待补}{测试集组织、事件检测与核验、人工评价工具及细则尚待确定；不表示完整benchmark已执行或发布。}

% % Paragraph role P09: 贡献入口
% Our contributions are as follows:
% \begin{itemize}
% \item We construct paired data for audiovisual backchannel generation during listening from Seamless conversations. Role-aware clip organization, media processing, and cue-assisted annotation link speaker audio, listener audiovisual media, and ordered reaction descriptions.
% \item We introduce a generation framework conditioned on current speaker audio and the listener's first frame. Direct audio conditioning supplies acoustic information, target feature supervision adapts the conditioning representations, and clip-level event supervision explicitly encourages feedback categories.
% \item We organize feedback evaluation around clip-level category agreement, contextual appropriateness and engagement, and audiovisual coordination, complementing media-quality evaluation while allowing reasonable variation in reaction timing.
% \end{itemize}
% \cnnote{贡献边界}{第三条保留为已确认的评价设计，不新增“首个”“首次”或“已发布benchmark”的主张。本文仅说明采用联合生成，不主张其已被证明优于音频—视频级联。}
